% Part: first-order-logic
% Chapter: syntax-and-semantics
% Section: substitution

\documentclass[../../../include/open-logic-section]{subfiles}

\begin{document}

\olfileid{fol}{syn}{sub}

\olsection{Substitusi}

\begin{defn}[Substitusi ing term]
Kita nemtokake $\Subst{s}{t}{x}$, yaiku asil saka \emph{nyubstitusikake}
$t$ kanggo saben muncule~$x$ ing $s$, kanthi rekursif:
\begin{enumerate}
\item \indcase{s}{c}{$\Subst{\indfrm}{t}{x}$ mung $s$.}

\item \indcase{s}{y}{$\Subst{\indfrm}{t}{x}$ uga mung~$s$,
  yen $y$ iku variabel lan $y \not\ident x$.}

\item \indcase{s}{x}{$\Subst{\indfrm}{t}{x}$ yaiku~$t$.}

\item\indcase{s}{\Atom{f}{t_1, \dots, t_n}}{$\Subst{\indfrmp}{t}{x}$
  yaiku $\Atom{f}{\Subst{t_1}{t}{x}, \dots, \Subst{t_n}{t}{x}}$.}
\end{enumerate}
\end{defn}

\begin{defn}
Term~$t$ iku \emph{!!{free for}} $x$ ing $!A$ yen ora ana muncule
bebas~$x$ ing $!A$ kang dumunung ing cakupan kuantor kang ngiket
sawijining variabel ing~$t$.
\end{defn}

\begin{ex} ~
\begin{enumerate}
\item $\Obj v_8$ iku bebas kanggo $\Obj v_1$ ing $\lexists[\Obj
  v_3]\Atom{\Obj A^2_4}{\Obj v_3,\Obj v_1}$

\item $\Obj f^2_1(\Obj v_1, \Obj v_2)$ \emph{ora} bebas kanggo $\Obj
  v_0$ ing $\lforall[\Obj v_2]\Atom{\Obj A^2_4}{\Obj v_0,\Obj v_2}$
\end{enumerate}
\end{ex}

\begin{defn}[Substitusi ing !!a{formula}]
Yen $!A$ iku !!a{formula}, $x$~iku !!a{variable}, lan $t$~iku term
kang !!{free for}~$x$ ing~$!A$, mula $\Subst{!A}{t}{x}$ yaiku asil saka
nyubstitusikake $t$ kanggo kabeh muncule bebas~$x$ ing~$!A$.
\begin{enumerate}
\tagitem{prvFalse}{\indcase{!A}{\lfalse}{$\Subst{\indfrm}{t}{x}$
    yaiku $\lfalse$.}}{}

\tagitem{prvTrue}{\indcase{!A}{\ltrue}{$\Subst{\indfrm}{t}{x}$
    yaiku $\ltrue$.}}{}

\item \indcase{!A}{\Atom{P}{t_1,\dots,
    t_n}}{$\Subst{\indfrm}{t}{x}$ yaiku $\Atom{P}{\Subst{t_1}{t}{x},
    \dots, \Subst{t_n}{t}{x}}$.}

\item \indcase{!A}{\eq[t_1][t_2]}{$\Subst{\indfrmp}{t}{x}$ yaiku
  $\Subst{t_1}{t}{x} = \Subst{t_2}{t}{x}$.}

\tagitem{prvNot}{\indcase{!A}{\lnot !B}{$\Subst{\indfrmp}{t}{x}$
    yaiku $\lnot \Subst{!B}{t}{x}$.}}{}

\tagitem{prvAnd}{\indcase{!A}{(!B \land
    !C)}{$\Subst{\indfrmp}{t}{x}$ yaiku $(\Subst{!B}{t}{x} \land
    \Subst{!C}{t}{x})$.}}{}

\tagitem{prvOr}{\indcase{!A}{(!B \lor
    !C)}{$\Subst{\indfrmp}{t}{x}$ yaiku $(\Subst{!B}{t}{x} \lor
    \Subst{!C}{t}{x})$.}}{}

\tagitem{prvIf}{\indcase{!A}{(!B \lif
    !C)}{$\Subst{\indfrmp}{t}{x}$ yaiku $(\Subst{!B}{t}{x} \lif
    \Subst{!C}{t}{x})$.}}{}

\tagitem{prvIff}{\indcase{!A}{(!B \liff
    !C)}{$\Subst{\indfrmp}{t}{x}$ yaiku $(\Subst{!B}{t}{x} \liff
    \Subst{!C}{t}{x})$.}}{}

\tagitem{prvAll}{
\indcase{!A}{\lforall[y][!B]}{$\Subst{\indfrmp}{t}{x}$
    yaiku $\lforall[y][\Subst{!B}{t}{x}]$, yen $y$ iku variabel
    saliyane $x$; yen ora, $\Subst{\indfrmp}{t}{x}$
    mung $\indfrm$.}}{}

\tagitem{prvEx}{
\indcase{!A}{\lexists[y][!B]}{$\Subst{\indfrmp}{t}{x}$
    yaiku $\lexists[y][\Subst{!B}{t}{x}]$, yen $y$ iku variabel
    saliyane $x$; yen ora, $\Subst{\indfrmp}{t}{x}$
    mung $\indfrm$.}}{}
\end{enumerate}
\end{defn}

\begin{explain}
Gatekna yen substitusi bisa tanpa pangaribawa: yen $x$ ora muncul ing
$!A$ babar pisan, mula $\Subst{!A}{t}{x}$ mung~$!A$.

Syarat manawa $t$ kudu !!{free for}~$x$ ing~$!A$ perlu kanggo
nyingkirake kasus kaya ing ngisor iki. Yen $!A \ident \lexists[y][x < y]$
lan $t \ident y$, mula $\Subst{!A}{t}{x}$ bakal dadi $\lexists[y][y <
  y]$. Ing kasus iki, variabel bebas $y$ ditangkep dening
kuantor $\lexists[y]$ nalika substitusi, lan iki ora dikarepake.
Contone, kita pengin yen $\lforall[x][!B]$ bener, banjur
$\Subst{!B}{t}{x}$ uga bener. Nanging gatekna
$\lforall[x][\lexists[y][x < y]]$ (ing kene $!B$ yaiku $\lexists[y][x <
  y]$). Iki ukara kang bener, upamane, tumrap wilangan asli:
kanggo saben wilangan~$x$ ana wilangan~$y$ kang luwih gedhe tinimbang
wilangan kasebut. Yen $y$ diidini dadi panggantos kanggo~$x$, kita bakal entuk
$\Subst{!B}{y}{x} \ident \lexists[y][y < y]$, kang luput. Kita
nyegah iki kanthi syarat supaya ora ana variabel bebas ing~$t$ kang
bakal dadi kaiket dening kuantor ing~$!A$.
\end{explain}

Kita kerep nggunakake pranatan ing ngisor iki supaya notasi ora
ruwet: yen $!A$ iku !!a{formula} kang bisa ngemot !!{variable}~$x$
kanthi bebas, kita uga nulis~$!A(x)$ kanggo nuduhake iki. Yen wis
cetha $!A$ lan~$x$ kang dimaksud, lan $t$ iku term (dianggep bebas
kanggo $x$ ing $!A(x)$), mula kita nulis $!A(t)$ minangka cekakan
kanggo $\Subst{!A}{t}{x}$. Dadi, upamane, kita bisa ngarani
$!A(t)$ sawijining instansi saka~$\lforall[x][!A(x)]$.
Tegese, yen $!A$~iku sembarang !!{formula}, $x$~!!a{variable},
lan $t$~term kang bebas kanggo~$x$ ing~$!A$, mula
$\Subst{!A}{t}{x}$ iku instansi saka~$\lforall[x][!A]$.

\end{document}
